\chapter{2009年四川卷（理）}

\section{单选题}

\begin{problem}
设集合 \(S = \{x \mid |x| < 5\}\)，\(T = \{x \mid x^2 + 4x - 21 < 0\}\)，则 \(S \cap T =\)（\quad）

\choices
    {\(\{x\mid -7 < x < -5\}\)}
    {\(\{x\mid 3 < x < 5\}\)}
    {\(\{x\mid -5 < x < 3\}\)}
    {\(\{x\mid -7 < x < 5\}\)}
\end{problem}

\begin{problem}
已知函数 \(f(x) = \begin{cases} a + \log_2 x & \text{当 } x \geqslant 2 \text{ 时} \\ \frac{x^2 - 4}{x - 2} & \text{当 } x < 2 \text{ 时} \end{cases}\) 在点 \(x = 2\) 处连续，则常数（\quad）

\choices
    {2}
    {3}
    {4}
    {5}
\end{problem}

\begin{problem}
复数 \(\frac{(1 + 2\i)^2}{3 - 4\i}\) 的值是（\quad）

\choices
    {\(-1\)}
    {1}
    {\(-\i\)}
    {\(\i\)}
\end{problem}

\begin{problem}
已知函数 \(f(x) = \sin \left(x - \frac{\pi}{2}\right)\)（\(x \in \R\)），下面结论错误的是（\quad）

\choices
    {函数 \( f(x) \) 的最小正周期为 \( 2\pi \)}
    {函数 \( f(x) \) 在区间 \(\left[0, \frac{\pi}{2}\right]\) 上是增函数}
    {函数 \( f(x) \) 的图象关于直线 \( x = 0 \) 对称}
    {函数 \( f(x) \) 是奇函数}
\end{problem}

\begin{problem}
如图，已知六棱锥 \(P - ABCDEF\) 的底面是正六边形，\(PA \perp\) 平面 \(ABC\)，\(PA = 2AB\)，则下列结论正确的是（\quad）

\bitmapfigure[max width=0.44\linewidth]{img/2009/sichuan_science/q5_fig1.png}

\choices
    {\(PB \perp AD\)}
    {平面 \(PAB \perp\) 平面 \(PBC\)}
    {直线 \(BC \parallel\) 平面 \(PAE\)}
    {直线 \(PD\) 与平面 \(ABC\) 所成的角为 \(45^{\circ}\)}
\end{problem}

\begin{problem}
已知 \(a\)，\(b\)，\(c\)，\(d\) 为实数，且 \(c > d\). 则“\(a > b\)”是“\(a - c > b - d\)”的（\quad）

\choices
    {充分而不必要条件}
    {必要而不充分条件}
    {充要条件}
    {既不充分也不必要条件}
\end{problem}

\begin{problem}
已知双曲线 \(\frac{x^2}{2} - \frac{y^2}{b^2} = 1\)（\(b > 0\)）的左，右焦点分别为 \(F_1\)，\(F_2\)，其一条渐近线方程为 \(y = x\)，点 \(P(\sqrt{3}, y_0)\) 在该双曲线上，则 \(\overrightarrow{PF_1} \cdot \overrightarrow{PF_2} =\)（\quad）

\choices
    {\(-12\)}
    {\(-2\)}
    {0}
    {4}
\end{problem}

\begin{problem}
如图，在半径为 3 的球面上有 \(A\)，\(B\)，\(C\) 三点，\(\angle ABC = 90^{\circ}\)，\(BA = BC\)，球心 \(O\) 到平面 \(ABC\) 的距离是 \(\frac{3\sqrt{2}}{2}\)，则 \(B\)，\(C\) 两点的球面距离是（\quad）

\bitmapfigure[max width=0.34\linewidth]{img/2009/sichuan_science/q8_fig1.png}

\choices
    {\(\frac{\pi}{3}\)}
    {\(\pi\)}
    {\(\frac{4\pi}{3}\)}
    {\(2\pi\)}
\end{problem}

\begin{problem}
已知直线 \(l_1: 4x - 3y + 6 = 0\) 和直线 \(l_2: x = -1\)，抛物线 \(y^2 = 4x\) 上一动点 \(P\) 到直线 \(l_1\) 和直线 \(l_2\) 的距离之和的最小值是（\quad）

\choices
    {2}
    {3}
    {\(\frac{11}{5}\)}
    {\(\frac{37}{16}\)}
\end{problem}

\begin{problem}
某企业生产甲，乙两种产品，已知生产每吨甲产品要用 \(A\) 原料 3 吨，\(B\) 原料 2 吨；生产每吨乙产品要用 \(A\) 原料 1 吨，\(B\) 原料 3 吨. 销售每吨甲产品可获得利润 5 万元. 每吨乙产品可获得利润 3 万元. 该企业在一个生产周期内消耗 \(A\) 原料不超过 13 吨，\(B\) 原料不超过 18 吨，那么该企业可获得最大利润是（\quad）

\choices
    {12 万元}
    {20 万元}
    {25 万元}
    {27 万元}
\end{problem}

\begin{problem}
3位男生和3位女生共6位同学站成一排，若男生甲不站两端，3位女生中有且只有两位女生相邻，则不同排法的种数是（\quad）

\choices
    {360}
    {188}
    {216}
    {96}
\end{problem}

\begin{problem}
已知函数 \( f(x) \) 是定义在实数集 \( \R \) 上的不恒为零的偶函数，且对任意实数 \( x \) 都有 \( xf(x + 1) = (1 + x)f(x) \)，则 \( f\left(f\left(\frac{5}{2}\right)\right) \) 的值是（\quad）

\choices
    {0}
    {\(\frac{1}{2}\)}
    {1}
    {\(\frac{5}{2}\)}
\end{problem}

\section{填空题}

\begin{problem}
\(\left(2x - \frac{1}{2x}\right)^6\) 的展开式的常数项是\fillinblank{}.（用数字作答）
\end{problem}

\begin{problem}
若 \(\odot O_1: x^{2} + y^{2} = 5\) 与 \(\odot O_2: (x - m)^{2} + y^{2} = 20\)（\(m \in \R\)）相交于 \(A\)，\(B\) 两点，且两圆在点 \(A\) 处的切线互相垂直，则线段 \(AB\) 的长度是\fillinblank{}.
\end{problem}

\begin{problem}
如图，已知正三棱柱 \(ABC - A_{1}B_{1}C_{1}\) 的各条棱长都相等，\(M\) 是侧棱 \(CC_{1}\) 的中点，则异面直线 \(AB_{1}\) 和 \(BM\) 所成的角的大小是\fillinblank{}.

\bitmapfigure[max width=0.42\linewidth]{img/2009/sichuan_science/q15_fig1.png}
\end{problem}

\begin{problem}
设 \(V\) 是已知平面 \(M\) 上所有向量的集合. 对于映射 \(f: V \to V\)，\(\bs{a} \in V\)，记 \(\bs{a}\) 的象为 \(f(\bs{a})\). 若映射 \(f: V \to V\) 满足：对所有 \(\bs{a}\)，\(\bs{b} \in V\) 及任意实数 \(\lambda\)，\(\mu\) 都有 \(f(\lambda\bs{a} + \mu\bs{b}) = \lambda f(\bs{a}) + \mu f(\bs{b})\)，则 \(f\) 称为平面 \(M\) 上的线性变换. 现有下列命题：

\begin{circlelist}
\item 设 \(f\) 是平面 \(M\) 上的线性变换，则 \(f(0) = 0\)；
\item 对 \(\bs{a} \in V\)，设 \(f(\bs{a}) = 2\bs{a}\)，则 \(f\) 是平面 \(M\) 上的线性变换；
\item 若 \(\bs{e}\) 是平面 \(M\) 上的单位向量，对 \(\bs{a} \in V\)，设 \(f(\bs{a}) = \bs{a} - \bs{e}\)，则 \(f\) 是平面 \(M\) 上的线性变换；
\item 设 \(f\) 是平面 \(M\) 上的线性变换，\(\bs{a}\)，\(\bs{b} \in V\)，若 \(\bs{a}\)，\(\bs{b}\) 共线，则 \(f(\bs{a})\)，\(f(\bs{b})\) 也共线.
\end{circlelist}

其中真命题是\fillinblank{}.（写出所有真命题的序号）
\end{problem}

\section{解答题}

\begin{problem}
在 \(\triangle ABC\) 中，\(A\)，\(B\) 为锐角，角 \(A\)，\(B\)，\(C\) 所对应的边分别为 \(a\)，\(b\)，\(c\)，且 \(\cos 2A = \frac{3}{5}\)，\(\sin B = \frac{\sqrt{10}}{10}\).

\begin{enumerate}
\item 求 \(A + B\) 的值；
\item 若 \(a - b = \sqrt{2} - 1\)，求 \(a\)，\(b\)，\(c\) 的值.
\end{enumerate}
\end{problem}

\begin{problem}
为振兴旅游业，四川省 2009 年面向国内发行总量为 2000 万张的熊猫优惠卡，向省外人士发行的是熊猫金卡（简称金卡），向省内人士发行的是熊猫银卡（简称银卡）. 某旅游公司组织了一个有 36 名游客的旅游团到四川名胜旅游，其中 \( \frac{3}{4} \) 是省外游客，其余是省内游客. 在省外游客中有 \( \frac{1}{3} \) 持金卡，在省内游客中有 \( \frac{2}{3} \) 持银卡.

\begin{enumerate}
\item 在该团中随机采访 3 名游客，求恰有 1 人持金卡且持银卡者少于 2 人的概率；
\item 在该团的省内游客中随机采访 3 名游客，设其中持银卡人数为随机变量 \(\xi\)，求 \(\xi\) 的分布列及数学期望 \(E(\xi)\).
\end{enumerate}
\end{problem}

\begin{problem}
如图，正方形 \(ABCD\) 所在平面与平面四边形 \(ABEF\) 所在平面互相垂直，\(\triangle ABE\) 是等腰直角三角形，\(AB = AE\)，\(FA = FE\)，\(\angle AEF = 45^{\circ}\).

\begin{enumerate}
\item 求证： \(EF \perp\) 平面 \(BCE\)；
\item 设线段 \(CD\) 的中点为 \(P\)，在直线 \(AE\) 上是否存在一点 \(M\)，使得 \(PM \parallel\) 平面 \(BCE\)？若存在，请指出点 \(M\) 的位置，并证明你的结论；若不存在，请说明理由；
\item 求二面角 \(F - BD - A\) 的大小.
\end{enumerate}

\bitmapfigure[max width=0.40\linewidth]{img/2009/sichuan_science/q19_fig1.png}

\end{problem}

\begin{problem}
已知椭圆 \(\frac{x^2}{a^2} + \frac{y^2}{b^2} = 1\)（\(a > b > 0\)）的左，右焦点分别为 \(F_1\)，\(F_2\)，离心率 \(e = \frac{\sqrt{2}}{2}\)，右准线方程为 \(x = 2\).

\begin{enumerate}
\item 求椭圆的标准方程；
\item 过点 \(F_{1}\) 的直线 \(l\) 与该椭圆交于 \(M\)，\(N\) 两点，且 \(\left|\overrightarrow{F_2M} + \overrightarrow{F_2N}\right| = \frac{2\sqrt{26}}{3}\)，求直线 \(l\) 的方程.
\end{enumerate}
\end{problem}

\begin{problem}
已知 \(a > 0\)，且 \(a \neq 1\) 函数 \(f(x) = \log_a(1 - a^x)\).

\begin{enumerate}
\item 求函数 \(f(x)\) 的定义域，并判断 \(f(x)\) 的单调性；
\item 若 \(n \in \mathbf{N}^*\)，求 \(\lim\limits_{n \to +\infty} \frac{a^{f(n)}}{a^n + a}\)；
\item 当 \(a = \e\)（\(\e\) 为自然对数的底数）时，设 \(h(x) = (1 - \e^{f(x)})(x^2 - m + 1)\). 若函数 \(h(x)\) 的极值存在，求实数 \(m\) 的取值范围以及函数 \(h(x)\) 的极值.
\end{enumerate}
\end{problem}

\begin{problem}
设数列 \(\{a_{n}\}\) 的前 \(n\) 项和为 \(S_{n}\)，对任意的正整数 \(n\)，都有 \(a_{n} = 5S_{n} + 1\) 成立，记 \(b_{n} = \frac{4 + a_{n}}{1 - a_{n}}\)（\(n \in \mathbf{N}^{*}\)）.

\begin{enumerate}
\item 求数列 \(\{b_{n}\}\) 的通项公式；
\item 记 \(c_{n} = b_{2n} - b_{2n-1}\)（\(n \in \mathbf{N}^{*}\)），设数列 \(\{c_{n}\}\) 的前 \(n\) 项和为 \(T_{n}\)，求证：对任意正整数 \(n\) 都有 \(T_{n} < \frac{3}{2}\)；
\item 设数列 \(\{b_{n}\}\) 的前 \(n\) 项和为 \(R_{n}\). 已知正实数 \(\lambda\) 满足：对任意正整数 \(n\)，\(R_{n} \leqslant \lambda n\) 恒成立，求 \(\lambda\) 的最小值.
\end{enumerate}
\end{problem}
